% GENERATED_BY: run_oc_core_monograph_translation_rework_dispatch_v1.py
% AI_ASSISTED_TRANSLATION_DRAFT: TRUE
% THEOREM_REVIEW_STATUS: PENDING
% THEOREM_REVIEW_MODE: FORMAL_AI_GATE__CODEX_CLI_CHATGPT
% THEOREM_REVIEW_REVIEWER_ID: 
% THEOREM_REVIEW_AUTHORITY_CLASS: 
% THEOREM_REVIEWED_AT: 
% THEOREM_REVIEW_DOSSIER_REF: 
% SOURCE_LANGUAGE: EN
% TARGET_LANGUAGE: RU
% SOURCE_EN_REF: content/04_results.tex
\section{Результаты и вытекающие следствия}
\label{sec:results}

В этом разделе суммируются ключевые структурные результаты, вытекающие из формального каркаса, введённого в разделе~\ref{sec:model}.
В отличие от более ранних черновиков Core, эта глава не является заглушкой: в ней изложены канонические следствия аксиом и операторов Ontology of Continua (OC), собранные по итогам более ранних внутренних прогонов и реорганизованные для монографии Core~1.3.

Все результаты этой главы предметно-независимы.
Они следуют исключительно из структурных определений континуумов, осей, порогов, потенциалов, потоков, границ, циклов и меры континуумности \(k(K,t)\).
Полные доказательства здесь не воспроизводятся; они вынесены в специальные статьи-расширения.
Основное внимание уделено формулировке ключевых теорем и выводов, задающих ограничения для всех континуумов в иерархии \(K_0\)–\(K_{10}\).

\subsection{Теорема 1: Монотонность размерности}

\textbf{Утверждение.}
Если континуум \(K_{x+1}\) возникает из \(K_x\), то
\[
    \dim(K_{x+1}) > \dim(K_x).
\]

\textbf{Обоснование (набросок).}
Размерностное возникновение определяется как появление новой оси \(A_{\mathrm{new}}\), несовместимой со всеми существующими осями из \(A(K_x)\); то есть её нельзя представить как их линейную или структурную комбинацию.
Размерностный порог \(\Theta_{\mathrm{dim}}(K_x)\) требует, чтобы \(A_{\mathrm{new}}\) была невырожденной и структурно необходимой.
Следовательно, любой возникающий континуум имеет строго большую размерность, чем его предшественник.

\textbf{Следствие 1.1 (Отсутствие деградационного упрощения).}
Не существует эволюции на уровне континуума, которая уменьшала бы размерность при сохранении существования:
если \(K(t)\) жив (\(k(K,t) > 0\)), то
\[
    \dim\big(K(t+dt)\big) \ge \dim\big(K(t)\big).
\]
Видимое уменьшение размерности соответствует смерти более высокоразмерного континуума и рождению другого, более низкоразмерного, а не непрерывному упрощению.

\subsection{Теорема 2: Невозможность спонтанного возникновения размерности}

\textbf{Утверждение.}
Ни один континуум не может увеличить свою размерность без превышения размерностного порога и без доступа к подходящей оси в своём пространстве вложения:
\[
    \dim(K_{x+1}) > \dim(K_x)
    \quad\Rightarrow\quad
    T(K_x,t) > \Theta_{\mathrm{dim}}(K_x)
    \ \wedge\
    A_{\mathrm{new}} \in A(M_x)\setminus A(K_x).
\]

\textbf{Обоснование (набросок).}
Новые оси представляют классы различий, несовместимые с уже существующими осями.
Такая несовместимость возникает только тогда, когда структурное напряжение \(T(K_x,t)\) относительно текущего порогового ландшафта не может быть разрешено в пределах существующих осей.
По определению \(\Theta_{\mathrm{dim}}\), ниже этого порога все различия можно спроецировать на линейную оболочку \(A(K_x)\); выше него проекция перестаёт работать, и требуется новая ось.
Новая ось уже должна быть доступна в пространстве вложения \(M_x\); в противном случае отсутствует структурное направление, вдоль которого континуум может расширяться.

\textbf{Вывод 2.1.}
Шум, локальные флуктуации или внутренние потоки, которые не поднимают структурное напряжение выше \(\Theta_{\mathrm{dim}}(K_x)\), не могут порождать новые размерности.

\subsection{Теорема 3: Смерть вследствие коллапса границы и пространства вложения}

\textbf{Утверждение.}
Континуум \(K\) умирает, когда коллапсирует его допустимое пространство состояний:
\[
    \Omega(K) = \emptyset.
\]

\textbf{Эквивалентные формулировки.}
Следующие условия эквивалентны и характеризуют смерть \(K\):
\begin{enumerate}
    \item поддерживающие циклы исчезают, \(C(K) = \emptyset\);
    \item поддерживающие потоки исчезают, \(J_{\mathrm{support}} = 0\), либо оказываются недостаточными для поддержания каких бы то ни было циклов;
    \item существует не зависящее от состояния нарушение порогов, \(\forall s:\ \exists\,k\ \text{такое, что } f_k(s) > 0\);
    \item мера континуумности коллапсирует, \(k(K,t) \to 0\);
    \item континуум покидает область, поддерживаемую его пространством вложения \(M\), то есть ни одно состояние не удовлетворяет одновременно внутренним порогам \(K\) и внешним ограничениям \(M\).
\end{enumerate}

\textbf{Следствие 3.1 (Смерть как выход из пространства вложения).}
Если ограничения \(M\) изменяются так, что никакую конфигурацию \(K\) больше нельзя вложить в \(M\) при соблюдении её порогов, то \(\Omega(K)\) становится пустым, и континуум умирает независимо от своей внутренней структуры.

\textbf{Следствие 3.2 (Необратимость смерти).}
Как только \(\Omega(K) = \emptyset\), никакой оператор \(E\), действующий в пределах того же уровня, и никакой оператор взаимодействия \(E_{\mathrm{int}}\), сохраняющий тождество \(K\), не может восстановить непустое \(\Omega(K)\).
Любое кажущееся ``восстановление'' соответствует рождению нового континуума \(K'\), а не воскрешению исходного \(K\).

\subsection{Теорема 4: Невозможность эволюции вне пространства вложения}

\textbf{Утверждение.}
Пусть \(K_x\) --- континуум, вложенный в \(M_x\).
Тогда для всех моментов времени \(t\),
\[
    A\big(K_x(t)\big) \subseteq A(M_x),
    \qquad
    \mathrm{span}\,A\big(K_x(t+dt)\big) \subseteq \mathrm{span}\,A(M_x).
\]

\textbf{Вывод 4.1.}
Все динамические процессы \(K_x\), включая размерностное рождение \(K_x \to K_{x+1}\), требуют, чтобы пространство вложения \(M_x\) уже содержало оси, вдоль которых протекает эволюция.
Континуум не может эволюционировать в направлениях, отсутствующих в его пространстве вложения.

\subsection{Теорема 5: Необходимость и достаточность совместимости с \texorpdfstring{$M$}{M}}

\textbf{Утверждение.}
Континуум \(K\) существует как живой континуум (то есть \(\Omega(K)\neq\emptyset\) и \(k(K,t)>0\)) тогда и только тогда, когда он совместим со своим пространством вложения \(M\):
\[
    K \text{ существует }
    \Longleftrightarrow
    \Omega(K) \neq \emptyset
    \ \wedge\
    A(K) \subseteq A(M)
    \ \wedge\
    \Theta(K) \text{ реализуема в } M.
\]

\textbf{Вывод 5.1.}
Даже если конфигурация математически корректно определена на уровне \(K\), она не соответствует живому континууму, если пространство вложения не может поддержать необходимые оси, пороги и потоки.
Несовместимые конфигурации структурно исключаются условием \(\Omega(K) = \emptyset\).

\subsection{Теорема 6: Структурное напряжение и фазовые переходы}

\textbf{Утверждение.}
Пусть \(T(K,t)\) обозначает структурное напряжение, понимаемое как функционал потенциалов, осей и их градиентов:
\[
    T(K,t) = T\big(P(t),A,\nabla P(t)\big).
\]
Тогда:
\begin{itemize}
    \item фазовый переход происходит, когда
          \(
              T(K,t) = \Theta_{\mathrm{crit}}(K),
          \)
    \item размерностный переход происходит, когда
          \(
              T(K,t) = \Theta_{\mathrm{dim}}(K),
          \)
    \item структурный коллапс происходит, когда
          \(
              T(K,t) = \Theta_{\mathrm{death}}(K).
          \)
\end{itemize}

\textbf{Следствие 6.1.}
Все качественные изменения континуума --- возникновение новых режимов, новых размерностей или коллапс --- управляются соотношением между структурным напряжением и пороговым ландшафтом.

\subsection{Теорема 7: Универсальный закон роста сложности}

\textbf{Утверждение.}
Для любого живого континуума \(K\),
\[
    S\big(K(t+dt)\big) \ge S\big(K(t)\big),
\]
где \(S\) --- это структурная мера сложности, определённая на пространстве состояний и его организации (например, объединяющая вклады числа осей, разнообразия состояний, богатства циклов и свойств пространства вложения).

\textbf{Вывод 7.1.}
Сложность строго возрастает всякий раз, когда:
\begin{itemize}
    \item появляются новые оси;
    \item формируются новые устойчивые циклы;
    \item расширяется допустимое пространство состояний \(\Omega(K)\);
    \item пространство вложения \(M\) приобретает новые оси, релевантные для \(K\).
\end{itemize}
Стагнирующая или постоянная сложность соответствует тонко сбалансированной динамике ниже критических порогов; однако такой режим в общем случае неустойчив (см. Теорему~\ref{thm:no-eternal-stabilisation}).

\subsection{Теорема 8: Невозможность вечной стабилизации}
\label{thm:no-eternal-stabilisation}

\textbf{Утверждение.}
Не существует нетривиального живого континуума \(K\), который навсегда остаётся в фиксированной точке своего пространства состояний, сохраняя при этом положительную континуумность.
Формально, не существует решения, для которого
\[
    \frac{dP}{dt} = 0, \quad \frac{dJ}{dt} = 0, \quad \frac{d\Theta}{dt} = 0, \quad \frac{dk}{dt} = 0
\]
для всех \(t\), за исключением тривиального случая, когда \(K\) структурно заморожен и отделён от своего пространства вложения.

\textbf{Обоснование (набросок).}
Сами пространства вложения \(M_x\) эволюционируют и расширяются (Теорема~10).
Внутренние и внешние потоки не могут одновременно быть тождественно нулевыми на всех временных масштабах, не приводя либо к:
\begin{itemize}
    \item коллапсу континуума (утрате поддерживающих потоков), либо
    \item структурным изменениям осей, порогов и циклов.
\end{itemize}
Следовательно, любой нетривиальный живой континуум подвержен либо эволюции своей структуры, либо конечной смерти.

\subsection{Теорема 9: Смерть как утрата циклов}

\textbf{Утверждение.}
Живой континуум \(K\) умирает, когда исчезает его максимальный структурно устойчивый циклический комплекс:
\[
    C_{\max}(K) = \emptyset.
\]

\textbf{Обоснование (набросок).}
Циклы представляют собой самоподдерживающиеся потоки, удерживающие континуум вдали от его границы \(\partial\Omega\).
Если такого цикла не существует, то:
\begin{itemize}
    \item поддерживающие потоки не могут образовать замкнутые петли;
    \item траектории либо приближаются к \(\partial\Omega\), либо нарушают пороги;
    \item континуумность \(k(K,t)\) убывает к нулю.
\end{itemize}
Следовательно, утрата максимального циклического комплекса совпадает с коллапсом \(\Omega(K)\).

\textbf{Следствие 9.1.}
Смерть можно структурно обнаружить по исчезновению всех циклов, удовлетворяющих минимальному условию устойчивости (строго положительному расстоянию до \(\partial\Omega\)).

\subsection{Теорема 10: Монотонный рост пространств вложения}

\textbf{Утверждение.}
Пространства вложения образуют монотонную последовательность:
\[
    M_0 \subset M_1 \subset \dots \subset M_x \subset \dots,
\]
и всякий раз, когда возникает новый континуум \(K_{x+1}\), соответствующее пространство вложения \(M_{x+1}\) строго расширяет \(M_x\) по набору доступных осей:
\[
    A(M_{x}) \subset A(M_{x+1}).
\]

\textbf{Вывод 10.1.}
Рост континуумов по размерности и сложности неотделим от роста их пространств вложения.
Новые континуумы не могут возникать без расширения структурных возможностей, закодированных в \(M\).

\subsection{Теорема 11: Несовместимость порогов и выразительности}

\textbf{Утверждение.}
Континуум \(K\) коллапсирует, когда его оси становятся недостаточными для выражения требуемых различий:
\[
    \dim\big(A(K)\big) < \dim\big(\mathrm{Differences}(K)\big),
\]
где \(\mathrm{Differences}(K)\) --- это эффективная размерность структурно релевантных различений, налагаемых пространством вложения и порогами.

\textbf{Вывод 11.1.}
Коллапс может наступить даже без нарушения энергетических или динамических ограничений, если превышена его репрезентационная ёмкость.
В частности, когнитивные, социальные или метатеоретические континуумы могут умирать исключительно из-за утраты выразительной адекватности своих осей.

\subsection{Теорема 12: Неполнота пространств вложения}

\textbf{Утверждение.}
Для любого конечного пространства вложения \(M_x\) существуют потенциальные континуумы \(K'\), которые не могут быть реализованы внутри \(M_x\), потому что их оси или пороги несовместимы с \(A(M_x)\) или с ограничениями \(M_x\).

\textbf{Вывод 12.1.}
Ни одно отдельное пространство вложения \(M_x\) не является универсальным для всех возможных континуумов.
Последовательность пространств вложения сама должна расширяться и разнообразиться, чтобы вмещать новые структурные режимы.

\subsection{Следствия и предметно-независимые выводы}

Собирая приведённые выше теоремы, получаем следующие предметно-независимые выводы:

\begin{itemize}
    \item Размерность для живых континуумов строго монотонна; непрерывная деградация размерности невозможна.
    \item Размерность не может возникать спонтанно; для этого требуется структурное напряжение выше \(\Theta_{\mathrm{dim}}\) и наличие подходящих осей в пространстве вложения.
    \item Смерть характеризуется коллапсом \(\Omega(K)\), исчезновением циклов и необратимостью этого коллапса.
    \item Эволюция ограничена пространствами вложения; континуумы не могут эволюционировать в направлениях, которые \(M\) не поддерживает.
    \item Сложность живых континуумов имеет тенденцию к росту, движимая новыми осями, новыми циклами и расширяющимися пространствами состояний.
    \item Вечная точная стабилизация структурно исключена для нетривиальных континуумов; они либо эволюционируют, либо умирают.
    \item Репрезентационные и выразительные пределы могут вызывать коллапс даже тогда, когда энергетические и динамические пределы ещё не достигнуты.
\end{itemize}

Эти выводы единообразно справедливы для всех уровней \(K_0\)–\(K_{10}\); их предметно-специфическое содержание возникает только из интерпретации осей, потенциалов, порогов и потоков на каждом уровне.

\subsection{Следствия для вертикальной иерархии}

В вертикальной иерархии \(K_0\)–\(K_{10}\) структурные теоремы проявляются следующим образом:

\begin{itemize}
    \item \(K_0\): никакой динамики, никакого рождения или смерти; только структурные условия различимости.
    \item \(K_1\): классическое фазоподобное поведение с непрерывными потоками и базовыми порогами устойчивости/критичности.
    \item \(K_2\): физические пороги, включая перколяцию, переходы BKT, пороги когерентности/конденсации и механизмы генерации массы.
    \item \(K_3\)–\(K_4\): химические и пребиотические пороги (замыкание RAF, формирование мембран, осмотические и кривизненные пороги, редокс- и pH-пороги).
    \item \(K_5\): электрические пороги возбуждения, спайкинга и потоков, управляемых градиентами, в ранних нейронных системах.
    \item \(K_6\): логические и репрезентационные пороги когнитивного связывания, предсказания, устойчивости памяти и когерентности внутренних моделей.
    \item \(K_7\)–\(K_8\): социальные и цивилизационные пороги, управляющие доверием, институциональной устойчивостью, системной жизнестойкостью и коллапсом.
    \item \(K_9\)–\(K_{10}\): пороги выразительности, когерентности и согласованности в пространстве теорий, парадигм, онтологий и метатеоретических структур.
\end{itemize}

Во всех случаях действуют одни и те же структурные принципы --- монотонность размерности, порогово управляемые возникновение и коллапс, зависимость от пространств вложения и рост сложности --- тогда как меняется лишь конкретная интерпретация переменных.

\subsection{Резюме}

Результаты, представленные в этой главе, образуют структурный каркас OC.
Они утверждают, что размерность монотонна, что возникающие размерности требуют порогово индуцированных фазовых переходов, что смерть характеризуется коллапсом допустимых состояний и необратима, что эволюция ограничена пространствами вложения, что сложность живых континуумов имеет тенденцию к росту, что циклы являются носителями структурной жизни и что выразительные пределы могут вызывать коллапс.
Эти принципы применимы ко всей иерархии континуумов и определяют ограничения, в рамках которых все континуумы могут существовать, эволюционировать или умирать.
